% Standalone TikZ diagram: minimal unity-feedback closed-loop system.
% Only the input R and output Y are labeled; intermediate signals
% (error, control effort) and the controller block are absorbed into a
% single forward-path TF G.
% Build:  pdflatex closed_loop_minimal.tex

% ---- Preamble ------------------------------------------------------------
% standalone: crops the output PDF tightly to the picture (no page chrome).
% border=4pt: leaves a small margin around the figure.
\documentclass[border=4pt]{standalone}
\usepackage{tikz}

% TikZ libraries used below:
%   positioning  -- enables `right=1.4cm of rin`-style relative placement
%   calc         -- enables coordinate arithmetic, e.g. ($(a)!0.5!(b)$)
%   arrows.meta  -- modern arrow-tip system; provides Latex[length=,width=]
\usetikzlibrary{positioning, calc, arrows.meta}

\begin{document}

% ---- Picture options -----------------------------------------------------
% Options inside [ ... ] apply to every node and edge in this tikzpicture.
\begin{tikzpicture}[
  % Default arrow tip for any `->` edge below.
  >={Latex[length=1.5mm, width=1.2mm]},
  thick,                                                          % default line weight

  % Re-usable named styles, referenced via `style=block` / `style=sum`
  % (or simply `block` / `sum` since pgf accepts the bare style name).
  block/.style = {draw, rectangle, minimum height=2.5em, minimum width=4.5em},
  sum/.style   = {draw, circle, inner sep=0pt, minimum size=1.2em},
]

  % ---- Nodes -------------------------------------------------------------
  % `coordinate` = an invisible point (no shape, no border, no label),
  % used as an anchor for arrow endpoints outside any visible block.
  \node[coordinate]                    (rin)   {};

  % `right=1.4cm of rin` (from the `positioning` library) places this
  % node 1.4 cm to the right of node `rin`'s east edge.
  % The summing junction is just an empty circle; the +/- polarity
  % labels are added separately below so they sit outside the circle.
  \node[sum,   right=1.4cm of rin]     (sum)   {};

  % The forward-path block.  Label is plain math `$G$`.
  \node[block, right=1.6cm of sum]     (G)     {$G$};

  % Anchor point for the output arrow's tip, just to the right of G.
  \node[coordinate, right=1.4cm of G]  (yout)  {};

  % ---- Sum polarity labels ----------------------------------------------
  % Small `+` and `-` text nodes anchored at the NW and SW corners of `sum`.
  % `anchor=south east at (sum.north west)` means: place the south-east
  % corner of this label at sum's north-west corner.  Net effect: the
  % label sits just outside the circle, hugging the upper-left.
  \node[font=\scriptsize, inner sep=0pt, anchor=south east] at (sum.north west) {$+$};
  \node[font=\scriptsize, inner sep=0pt, anchor=north east] at (sum.south west) {$-$};

  % ---- Edges (signal arrows) --------------------------------------------
  % `--` draws a straight line between the two endpoints.
  % `->` uses the default arrow tip configured in the picture options.
  % `node[above]{...}` places a text label centered above the line.
  \draw[->] (rin) -- node[above]{$R$} (sum);   % reference R into summing junction
  \draw[->] (sum) -- (G);                      % error signal (unlabeled here)
  \draw[->] (G)   -- node[above]{$Y$} (yout);  % plant output Y

  % ---- Feedback path ----------------------------------------------------
  % `calc` syntax: ($A!t!B$) is the point at fractional distance t (0..1)
  % along the segment from A to B.  Here t=0.5 → `fbk` is the midpoint
  % between G's east edge and `yout`.  This is where the feedback tap
  % branches off the output line.
  \coordinate (fbk) at ($(G.east)!0.5!(yout)$);

  % Feedback wire: starting at `fbk`,
  %   `-- ++(0,-1.4)`  draws a straight segment by a *relative* offset of
  %                    (0, -1.4 cm) — i.e. straight down 1.4 cm.
  %   `-| (sum.south)` routes the rest orthogonally: first horizontal,
  %                    then vertical, ending at the south side of `sum`.
  \draw[->] (fbk) -- ++(0,-1.4) -| (sum.south);

\end{tikzpicture}
\end{document}
